\section{Orientation and the 2D Cross Product}

One of the most important predicates in computational geometry is the
orientation test. Given three points
\[
A=(x_A,y_A),\qquad B=(x_B,y_B),\qquad C=(x_C,y_C),
\]
we want to determine whether traveling from $A$ to $B$ and then toward $C$
represents a left turn, right turn, or no turn.

Define
\[
\overrightarrow{AB}=B-A, \qquad \overrightarrow{AC}=C-A.
\]
The two-dimensional cross product is
\[
\overrightarrow{AB}\times\overrightarrow{AC}
=(x_B-x_A)(y_C-y_A)-(y_B-y_A)(x_C-x_A).
\]

We therefore define
\[
\operatorname{orientation}(A,B,C)
=(x_B-x_A)(y_C-y_A)-(y_B-y_A)(x_C-x_A).
\]

The sign gives the geometric relationship:
\[
\operatorname{orientation}(A,B,C)
\begin{cases}
>0, & \text{counterclockwise / left turn},\\
<0, & \text{clockwise / right turn},\\
=0, & \text{collinear}.
\end{cases}
\]

\subsection{Geometric Interpretation}

The magnitude of the cross product is the area of the parallelogram formed
by $\overrightarrow{AB}$ and $\overrightarrow{AC}$. Therefore the area of
triangle $ABC$ is
\[
\operatorname{Area}(ABC)
=\frac{1}{2}\left|\operatorname{orientation}(A,B,C)\right|.
\]

If the orientation is zero, the triangle has zero area, meaning that the
three points are collinear.

\subsection{Floating-Point Considerations}

When coordinates are represented using floating-point values, exact
comparison with zero can be unreliable. Instead, a small tolerance
$\varepsilon$ is used:
\[
|x|<\varepsilon.
\]

In this project,
\[
\varepsilon=10^{-9}.
\]

Values sufficiently close to zero are therefore treated as zero when
checking collinearity.
